\documentclass{apcalc}
\title{Unit 3: Differentiation: Composite, Implicit, and Inverse Functions}
\author{Paul Wang}
\date{November 2026}
\volume{Study Guide 3 of 10 (topics 3.1--3.6)}
\shorttitle{Unit 3 --- Study Guide}
\begin{document}
\maketitle

\begin{abstract}
Unit 3 completes the toolkit: the chain rule, implicit differentiation, derivatives of inverse functions (including the inverse trig functions), and higher-order derivatives. It carries 5--10\% of the multiple-choice section, and the chain rule is in nearly every derivative you will take for the rest of the course. Free-response questions use these tools on implicitly defined curves (``find $\frac{dy}{dx}$; find where the tangent is horizontal'') and on functions given by tables. This guide states the unit in eight rules (R3.1--R3.8) across four sections. The \emph{Unit 3 Practice Set} uses the same section numbers.
\end{abstract}
\keywords{chain rule, composite function, implicit differentiation, horizontal and vertical tangents, inverse function, inverse trigonometric functions, second derivative.}

\section*{Router}
\begin{router}
a function inside a function: $f(g(x))$, $\sin(3x^2)$, $e^{-x/2}$ & chain rule: outside derivative at the inside, times inside derivative & R3.1 \\
$h(x)=f(g(x))$ with $f$, $g$ given by a table & $h'(a)=f'(g(a))\,g'(a)$: find $g(a)$ first & R3.2 \\
an equation in $x$ and $y$ not solved for $y$ & differentiate both sides, $\frac{d}{dx}y^n=ny^{n-1}\frac{dy}{dx}$ & R3.3 \\
horizontal / vertical tangent on an implicit curve & numerator $=0$ / denominator $=0$ of $\frac{dy}{dx}$, and the point is on the curve & R3.4 \\
$(f^{-1})'(b)$ & $\dfrac{1}{f'(a)}$ where $f(a)=b$ & R3.5 \\
$\arcsin$, $\arccos$, $\arctan$ & table of inverse trig derivatives & R3.6 \\
$\frac{d^2y}{dx^2}$ for an implicit curve & differentiate $\frac{dy}{dx}$ again, then substitute $\frac{dy}{dx}$ & R3.8 \\
\end{router}

\section{The chain rule (3.1)}

\begin{rulebox}{R3.1}{Chain rule}
$\dfrac{d}{dx}f(g(x))=f'(g(x))\,g'(x)$, or $\dfrac{dy}{dx}=\dfrac{dy}{du}\dfrac{du}{dx}$. Work from the outside in, one layer at a time:
$\frac{d}{dx}\sin^3(2x)=3\sin^2(2x)\cdot\cos(2x)\cdot2$. Common forms: $\frac{d}{dx}e^{u}=e^{u}u'$, $\frac{d}{dx}\ln u=\frac{u'}{u}$, $\frac{d}{dx}u^n=nu^{n-1}u'$.
\end{rulebox}

\begin{rulebox}{R3.2}{The chain rule with tables and graphs}
For $h(x)=f(g(x))$: $h'(a)=f'\big(g(a)\big)\cdot g'(a)$. Evaluate $g(a)$ first and feed that \emph{output} into $f'$. For $k(x)=g(x)^2$ use $k'(a)=2g(a)g'(a)$.
\end{rulebox}

\begin{example}{a composite from a table}
\begin{center}\small\begin{tabular}{|c|c|c|c|c|}\hline $x$ & $f(x)$ & $f'(x)$ & $g(x)$ & $g'(x)$\\\hline 1 & 3 & $-2$ & 4 & 5\\\hline 4 & 0 & 7 & 1 & $-3$\\\hline\end{tabular}\end{center}
$h(x)=f(g(x))$: $h'(1)=f'(g(1))\,g'(1)=f'(4)\cdot5=7\cdot5=35$. \answer{35.} (Not $f'(1)\,g'(1)=-10$.)
\end{example}

\section{Implicit differentiation (3.2)}

\begin{rulebox}{R3.3}{Procedure}
Differentiate both sides with respect to $x$, treating $y$ as a function of $x$: every $y$-term picks up a factor $\frac{dy}{dx}$ (chain rule), and products like $xy$ need the product rule. Collect the $\frac{dy}{dx}$ terms on one side, factor, and divide.
\end{rulebox}

\begin{example}{a tangent line to an implicit curve}
$x^2+xy+y^3=7$ at $(2,1)$. Differentiate: $2x+y+x\frac{dy}{dx}+3y^2\frac{dy}{dx}=0$, so $\dfrac{dy}{dx}=-\dfrac{2x+y}{x+3y^2}$. At $(2,1)$: $-\frac{5}{5}=-1$. \answer{$y=1-(x-2)$.}
\end{example}

\begin{rulebox}{R3.4}{Horizontal and vertical tangents}
Write $\frac{dy}{dx}=\frac{N}{D}$. Horizontal tangent: $N=0$ and $D\neq0$. Vertical tangent: $D=0$ and $N\neq0$. In both cases solve together with the original equation: the point must lie on the curve.
\end{rulebox}

\begin{frqtip}
``Is there a point on the curve where the tangent is horizontal?'' Set $N=0$, substitute into the curve's equation, and show either a real solution or a contradiction. A common pattern: $N=0$ forces $y=-2x$, and substituting gives an equation with no real solution, so the answer is no.
\end{frqtip}

\section{Derivatives of inverse functions (3.3--3.4)}

\begin{rulebox}{R3.5}{Inverse function rule}
If $f$ is differentiable and one-to-one with $f(a)=b$ and $f'(a)\neq0$, then
\[(f^{-1})'(b)=\frac{1}{f'(a)}=\frac{1}{f'\big(f^{-1}(b)\big)}.\]
The slopes at matching points $(a,b)$ on $f$ and $(b,a)$ on $f^{-1}$ are reciprocals. Find $a$ by solving $f(a)=b$ (often by inspection or from a table).
\end{rulebox}

\begin{example}{inverse from a formula}
$f(x)=x^3+2x+1$, $g=f^{-1}$. Find $g'(4)$. Solve $f(a)=4$: $a=1$ works ($1+2+1$). $f'(x)=3x^2+2$, $f'(1)=5$. \answer{$g'(4)=\tfrac15$.}
\end{example}

\begin{rulebox}{R3.6}{Inverse trig derivatives}
\[\frac{d}{dx}\arcsin x=\frac{1}{\sqrt{1-x^2}},\qquad \frac{d}{dx}\arccos x=-\frac{1}{\sqrt{1-x^2}},\qquad \frac{d}{dx}\arctan x=\frac{1}{1+x^2}.\]
With the chain rule: $\frac{d}{dx}\arctan(u)=\frac{u'}{1+u^2}$ and $\frac{d}{dx}\arcsin(u)=\frac{u'}{\sqrt{1-u^2}}$. (Also $\frac{d}{dx}\operatorname{arcsec}x=\frac{1}{|x|\sqrt{x^2-1}}$, rarely tested.)
\end{rulebox}

\begin{trap}
$\sin^{-1}x$ means $\arcsin x$, not $\frac{1}{\sin x}$. And $\frac{d}{dx}\arctan(2x)=\frac{2}{1+4x^2}$: square the whole inside, $u^2=(2x)^2$.
\end{trap}

\section{Selecting procedures and higher-order derivatives (3.5--3.6)}

\begin{rulebox}{R3.7}{Choosing the rule}
Ask what the \emph{last} operation is: a product, a quotient, a composition, or an inverse. Simplify before differentiating when you can ($\ln\frac{x^2}{\sqrt{x+1}}=2\ln x-\frac12\ln(x+1)$; $\frac{x^2+1}{x}=x+x^{-1}$).
\end{rulebox}

\begin{rulebox}{R3.8}{Higher-order and implicit second derivatives}
$f''=(f')'$; in Leibniz notation $\frac{d^2y}{dx^2}$. For an implicit curve, differentiate $\frac{dy}{dx}$ (quotient rule, with $y$ still a function of $x$), then substitute the expression for $\frac{dy}{dx}$, and at a point plug in numbers as early as possible.
\end{rulebox}

\begin{example}{implicit second derivative}
$x^2+y^2=25$: $\frac{dy}{dx}=-\frac{x}{y}$. Then $\frac{d^2y}{dx^2}=-\frac{y-x\frac{dy}{dx}}{y^2}=-\frac{y+\frac{x^2}{y}}{y^2}=-\frac{x^2+y^2}{y^3}=-\frac{25}{y^3}$. At $(3,4)$: $-\frac{25}{64}$. \answer{$-\tfrac{25}{64}$.}
\end{example}

\section{Common errors}
\begin{itemize}
\item Forgetting the inside derivative: $\frac{d}{dx}e^{3x}=e^{3x}$ (R3.1).
\item Evaluating $f'(a)$ instead of $f'(g(a))$ in a table composite (R3.2).
\item Writing $\frac{d}{dx}y^2=2y$ without $\frac{dy}{dx}$ (R3.3).
\item Using $f'(b)$ instead of $f'(a)$ for $(f^{-1})'(b)$ (R3.5).
\item Forgetting the minus sign on $\arccos$ (R3.6).
\end{itemize}

\end{document}
